\documentclass[10pt,a4paper]{amsart}
\usepackage[margin=19mm]{geometry}
\usepackage{amsmath,amssymb,mathtools}
\usepackage[scr=rsfs,cal=euler]{mathalfa}
\usepackage{xfrac}
\usepackage[unicode,hidelinks,bookmarks=false]{hyperref}
\newcommand{\ie}{i.e.\ }
\newcommand{\cf}{cf.\ }
\newcommand{\resp}{respectively\ }
\newcommand{\Subsection}{\subsection}
\providecommand{\nobreakdash}{\nobreak\hbox{-}}
\setlength{\emergencystretch}{2em}
\multlinegap0pt
\usepackage{bm}
\usepackage{bbm}
\usepackage{dsfont}
\usepackage{xspace}
\usepackage{cancel}
\usepackage{mathtools}
\usepackage[shortlabels]{enumitem}
\usepackage{amsthm}
\makeatletter
\@ifundefined{theorem}{\newtheorem{theorem}{Theorem}}{}
\@ifundefined{claim}{\newtheorem{claim}[theorem]{Claim}}{}
\@ifundefined{lemma}{\newtheorem{lemma}[theorem]{Lemma}}{}
\makeatother
\newcommand{\La}{\lambda}
\newcommand{\Sint}{\mathbf{S}}
\newcommand{\cone}{\operatorname{cone}}
\title[The number $4/9$ is a non-jump for $3$-graphs]{The number $4/9$ is a non-jump for $3$-graphs}
\author{Xizhi Liu, Dhruv Mubayi}
\date{Source: arXiv:2605.13567v1 (CC BY 4.0)}
\begin{document}
\maketitle

\noindent\emph{Source and licence.} The number $4/9$ is a non-jump for $3$-graphs. Version arXiv:2605.13567v1 (CC BY 4.0), distributed under the Creative Commons Attribution 4.0 International licence, as recorded in the arXiv licence metadata for this exact version and shown on the abstract page https://arxiv.org/abs/2605.13567v1. Full rights notice: licenses/arxiv-2605.13567-CC-BY-4.0.txt.\par\smallskip

\noindent\textit{Editorial scope.} Counted unit: the Lemma whose source label is \texttt{lem:code} in the article (source0.tex lines 532--539, source label \texttt{lem:code}), delivered together with the article's own complete original proof of it (source0.tex lines 565--607, one \texttt{proof} environment of 43 lines). No mathematics is written, repaired or re-derived: statement and proof are the article's own text line for line. The proof is detached from the statement in the source: the article states the result at lines 532--539 and prints its proof at lines 565--607, and the proof environment's own header names the statement, so the association is the article's own. Required context retained uncounted: thm:main (lines 94--96); lem:FR-criterion (lines 242--244); thm:local-cone (lines 273--275). Every definition, display and label of those lines is retained unchanged. Omitted from this package and not redistributed: 1--93, 97--241, 245--272, 276--531, 540--564, 608--845. Nothing inside a retained excerpt is omitted or altered: every retained body is delivered exactly as the article's lines are written, so no omission receipt of any kind accompanies this package. Citations: the retained excerpts carry no citation command at all, so no bibliography entry of the article is transcribed and no reference is redistributed. Reference adaptation: none was needed and none was made; every reference inside the delivered unit resolves inside it, so no unresolved reference or citation is produced. Printed slips kept literal: no printed word, symbol, hypothesis or number of the counted statement or of its proof was altered. Layout and class: the article's own class is \texttt{article} with packages including geometry, fontenc, inputenc, amsmath, amssymb, amsthm, mathtools, enumitem, microtype, tikz, float, xcolor, hyperref; the wrapper is the lane's pinned \texttt{amsart} editorial wrapper at 10pt on A4, which restates the theorem-like environments the retained text needs and adds the article's own macro definitions that the retained text needs (\texttt{\textbackslash La}, \texttt{\textbackslash Sint}, \texttt{\textbackslash cone}), with extra wrapper packages (bm, bbm, dsfont, xspace, cancel, mathtools, enumitem). The delivered numbering is the unit's own. Assets: no retained span contains a figure environment, a graphics-inclusion command, a PDF-page inclusion, a table or a TikZ picture; no image file is copied into this package, the package declares no asset dependency and no image file is redistributed with it. Licence: version arXiv:2605.13567v1 of this article is distributed under the Creative Commons Attribution 4.0 International licence, as recorded in the arXiv licence metadata for this exact version and shown on the abstract page https://arxiv.org/abs/2605.13567v1. A full rights notice is delivered at licenses/arxiv-2605.13567-CC-BY-4.0.txt. Characters: every retained excerpt is pure ASCII, so the delivered engine, XeLaTeX with its default Latin Modern text font, reports no missing character.
\begin{theorem}\label{thm:main}
    The number $4/9$ is a non-jump for $3$-uniform hypergraphs.
\end{theorem}

\begin{lemma}\label{lem:FR-criterion}
Let $\alpha\in[0,1)$. Suppose that for every integer $m$ there exists a finite $3$-graph $G_m$ such that $\La(G_m)>\alpha$, but $\La(H)\le \alpha$ for every subgraph $H\subseteq G_m$ with $v(H)\le m$. Then $\alpha$ is a non-jump.
\end{lemma}

\begin{theorem}\label{thm:local-cone}
Let $Q$ be a sparse $3$-graph with $\Delta_2(Q)\le2$. Then $\La(\cone(Q))\le 4/9$.
\end{theorem}

\begin{lemma}\label{lem:code}
For every integer $m$ there are arbitrarily large integers $t$ and a $3$-graph $\Sint$ on $t$ vertices such that
\begin{enumerate}[label=(\roman*)]
\item $\Delta_2(\Sint)\le2$;
\item $|\Sint|=t(t-1)/3$;
\item every induced subgraph of $\Sint$ on at most $m$ vertices is sparse.
\end{enumerate}
\end{lemma}

\begin{proof}[Proof of Theorem~\ref{thm:main}]
Fix an integer $m$. Let $\Sint$ be given by Lemma~\ref{lem:code} on a vertex set $B$ of size $t$, where $t$ is large. Define $G \coloneqq \cone(\Sint)$.

\begin{claim}\label{lem:global-lag}
For all sufficiently large $t$, we have $\La(G)>\frac49$.
\end{claim}

\begin{proof}
Let $x_{v_0} \coloneqq \frac13$ and $x_i \coloneqq \frac{2}{3t}$ for $i \in B$. The $v_0BB$-edges contribute
\[
        6\cdot\frac13\binom t2\left(\frac{2}{3t}\right)^2
        =\frac49\left(1-\frac1t\right).
\]
The edges inside $B$ contribute, using $|\Sint|=t(t-1)/3$,
\[
        6|\Sint|\left(\frac{2}{3t}\right)^3
        =6\cdot\frac{t(t-1)}3\cdot\frac{8}{27t^3}
        =\frac{16(t-1)}{27t^2}.
\]
Therefore
\[
        \La(G)
        \ge
        \frac49\left(1-\frac1t\right)+\frac{16(t-1)}{27t^2}
        =\frac49+\frac4{27t}-\frac{16}{27t^2}.
\]
This is greater than $4/9$ whenever $t>4$.
\end{proof}


\begin{claim}\label{lem:small-subgraphs}
Every subgraph $H\subseteq G$ with $v(H)\le m$ satisfies $\La(H)\le \frac49$.
\end{claim}

\begin{proof}
Let $U \coloneqq V(H)\cap B$. Then $|U|\le m$. By Lemma~\ref{lem:code}, the induced subgraph $\Sint[U]$ is sparse and satisfies $\Delta_2(\Sint[U])\le2$.

If $v_0\in V(H)$, then $H\subseteq \cone(\Sint[U])$. If $v_0\notin V(H)$, then $H\subseteq \Sint[U]$, and after adding an unused apex vertex we again have $H\subseteq \cone(\Sint[U])$. By Theorem~\ref{thm:local-cone}, $\La(\cone(\Sint[U]))\le\frac49$. Since normalized Lagrangian is monotone under adding edges and isolated vertices, it follows that $\La(H)\le4/9$.
\end{proof}


For every integer $m$, Claims~\ref{lem:global-lag} and~\ref{lem:small-subgraphs} give a finite $3$-graph $G$ such that $\La(G)>\frac49$ and $\La(H)\le\frac49$ for every subgraph $H\subseteq G$ with $v(H)\le m$. Lemma~\ref{lem:FR-criterion}, applied with $\alpha=4/9$, shows that $4/9$ is a non-jump.
\end{proof}

\end{document}
